% Arquivo de configuracao do documento
% Define os dados do autor e do trabalho

% Nome do autor, exceto o sobrenome
% Por exemplo, para Edgar Frank Codd, ficaria:
% \autorNomeA{Edgar Frank}
\autorNomeA{Kelly dos Reis} % Este eh o primeiro autor

% Sobrenome do autor
% Por exemplo, para Edgar Frank Codd, ficaria:
% \autorSobrenomeA{Codd}
\autorSobrenomeA{Leite}

% Nome do segundo autor (se houver), exceto o sobrenome
% Seguir o modelo anterior para o autor A
% Se nao houver segundo autor, deixe o comando abaixo vazio,
% isto eh, sem nada entre as chaves, inclusive espacos em branco: \autorNomeB{}
\autorNomeB{}

% Sobrenome do segundo autor
% Seguir o modelo anterior para o autor A
% Se nao houver segundo autor, deixe o comando abaixo vazio.
\autorSobrenomeB{}

% Contato do(s) autor(es): informar o email
\emailAutorA{kelly.l.reis@gmail.com}
\emailAutorB{}

% O titulo do trabalho
\titulo{Plataforma Web Integrada para Fórum Acadêmico e Voluntariado}

% Dados do orientador
% Nome do professor orientador, exceto sobrenome
% Por exemplo, para Luiz Olmes Carvalho, ficaria:
%\orientadorNome{Luiz Olmes}
\orientadorNome{Bruno Guazzelli}

% Sobrenome do orientador
% Por exemplo, para Luiz olmes Carvalho, ficaria:
% \orientadorSobrenome{Carvalho}
\orientadorSobrenome{Batista}

% Dados do coorientador
% Nome do professor coorientador, exceto sobrenome
% Por exemplo, para Rafael de Magalhaes Dias Frinhani, ficaria:
% \coorientadorNome{Rafael de Magalhaes Dias}
% Se o trabalho nao possui coorientador, deixe o comando abaixo vazio,
% isto eh, sem nada entre as chaves, inclusive espacos em branco: \coorientadorNome{}
\coorientadorNome{}

% Sobrenome do orientador
% Por exemplo, para Rafael de Magalhaes Dias Frinhani, ficaria:
% \coorientadorSobrenome{Frinhani}
\coorientadorSobrenome{}

% Definicao de textos: preencha adequadamente
% Opcoes: Orientador, Orientadora, Coorientador, Coorientadora
\txtOrientador{Orientador}
\txtCoorientador{Coorientador}

% Ano de publicacao do trabalho
\ano{2026}

% Mes da publicacao do trabalho
% Nao se trata do mes da defesa, mas sim do mes em que foi realizada a entrega da versao final.
\mes{Julho}

% Local: nao altere
\local{Itajubá}

% Palavras chave: incluidas no fim do abstact e na ficha catalografica.
% Minimo de 3 palavras e maximo de 5 palavras.
% Se nao possuir as quarta e quinta palavras, deixe o campo vazio, sem nada entre as chaves, inclusive espacos em branco. Ex: \palavrasChaveD{}
\palavraChaveA{Fórum Acadêmico}
\palavraChaveB{ONGs Voluntárias}
\palavraChaveC{Plataforma Web}
\palavraChaveD{Engajamento Comunitário}
\palavraChaveD{Integração Universitária}

% Nome dos membros da banca examinadora: maximo de 5 membros
% Se a banca possuir menos membros, deixe os comandos vazios,
% isto eh, sem nada entre as chaves, inclusive espacos em branco. Ex: \membroBanca5{}
\membroBancaA{} % Normalmente, o primeiro membro da banca eh o orientador
\membroBancaB{}
\membroBancaC{}
\membroBancaD{}
\membroBancaE{}

% Data da defesa: para fins de preenchimento da folha de aprovacao
% Coloque o mes por extenso, como no exemplo abaixo
\dataDefesa{}

% Apos a defesa, adicione o scan do PDF da folha de aprovacao assinada ao seu projeto.
% Use o comando abaixo para indicar o caminho do arquivo. Ex: \pdfAprovacao{./figs/folha.pdf}
\pdfAprovacao{} % Para a versao de defesa, deixe este comando vazio, sem nada entre as chaves

% Pacotes ja carregados:
% algorithm2e, amsmath, amssymb, amsthm, array
% babel, caption, xcolor, enumitem, eso-pic, float
% ifthen, fancybox, fancyhdr, fontenc, geometry, hyperref
% graphicx, hyphenat, lastpage, lipsum, listings
% longtable, lscape, multirow, natbib, pdfpages
% pxfonts, rotating, setspace, subfigure, textpos

% Se necessario, adicione pacotes aqui.
% Use o comando \usepackage{}
% \usepackage{nome-do-pacote}
\usepackage{float}